\documentclass[11pt,reqno]{amsart}
\usepackage[a4paper,margin=28mm]{geometry}
\usepackage{amsmath,amssymb,mathtools,mathrsfs}
\usepackage{enumitem,microtype,xurl}
\usepackage[colorlinks=true,linkcolor=blue,citecolor=blue,urlcolor=blue]{hyperref}
\numberwithin{equation}{section}
\setlist[enumerate]{label=\textup{(\roman*)},leftmargin=2.2em}
\newtheorem{theorem}{Theorem}[section]
\newtheorem{proposition}[theorem]{Proposition}
\newtheorem{lemma}[theorem]{Lemma}
\newtheorem{corollary}[theorem]{Corollary}
\theoremstyle{definition}
\newtheorem{definition}[theorem]{Definition}
\newtheorem{example}[theorem]{Example}
\theoremstyle{remark}
\newtheorem{remark}[theorem]{Remark}
\newcommand{\Nzero}{\mathbb N_0}
\newcommand{\K}{\mathcal K}
\newcommand{\Ssafe}{\mathscr S}
\newcommand{\Patt}{\mathsf P}
\newcommand{\hpat}{h^{*}}
\newcommand{\doi}[1]{\href{https://doi.org/#1}{doi:\nolinkurl{#1}}}

\title[Sparse sampling and Borel refinements]{Sparse Sampling and Borel Refinements of Closed Covers}
\author{Ziqing Ding}
\address{School of Mathematical Sciences, Nanjing Normal University,
Nanjing 210023, China}
\date{October 3, 2026}
\subjclass[2020]{37B10, 37B40, 93C55}
\keywords{sequence entropy, Borel partition, closed cover, symbolic dynamics,
safe feedback, maximal pattern entropy}

\begin{document}
\begin{abstract}
Let $T$ be conjugate to a positive power of a mixing one-sided shift of
finite type, and let $\K$ be a finite closed cover of its state space.
We prove that, along every increasing sequence whose successive gaps tend
to infinity, the minimum combinatorial sequence entropy over all finite
Borel partitions refining $\K$ is $\log N(\K)$, where $N(\K)$ is the
minimum size of a subcover. The same partition attains this value for all
such sequences and minimizes maximal pattern entropy. The sampling sequence
is fixed before the Borel partition is chosen. The proof uses the
nonmeager atoms of a partition, openness of the shift, and finite-type
concatenation. Applied to safe control blocks with a common successor, the
result gives an exact fixed-period minimax formula over all legal Borel
feedbacks, including arbitrary Borel mixtures and atom refinements.
A six-control system with two possible successors at each state has
minimum safe-cover size three but full-feedback minimax value $\log 2$.
Thus the cover formula need not survive feedback-dependent successors,
even with clopen safety domains and continuous open branches.
\end{abstract}
\maketitle

\section{Introduction}

For a finite partition observed along an orbit, one can count the names
seen at a prescribed sparse set of times. Two choices then arise: the
sampling sequence and the observation partition. When observations must
refine a given cover, it is natural to ask whether a single sequence,
chosen before the partition, detects the least possible complexity.
For Borel partitions this is not a finite-window question: the atoms need
not contain open sets, and the partition may depend on the entire state.

This paper determines that optimization for finite closed covers of
mixing shifts of finite type. If $r$ is the minimum size of a subcover,
then every sequence with gaps tending to infinity gives a lower bound
$\log r$ for \emph{every} finite Borel refinement. An ordered difference
partition of a minimum subcover attains the bound. In particular, the
sequence $(0,1,4,9,\ldots)$ can be specified in advance. The theorem is
about a common minimum threshold; it does not assert that this sequence
maximizes the entropy of each individual partition.

The distinction between a cover and its refining partitions is essential.
Romagnoli \cite{Romagnoli2007} computed, on the two-sided full shift over
$q>2$ symbols, a clopen cover with consecutive-time cover entropy
$\log(q/(q-1))$ and minimum consecutive-time entropy $\log 2$ among its
clopen refinements. The cover excludes the current symbol. Our result
allows all Borel refinements, uses prescribed sparse times, and applies
to arbitrary finite closed covers of mixing finite-type shifts. We spell
out the corresponding sparse cover--partition gap in
Corollary~\ref{cor:covergap}. This does not settle the consecutive-time
Borel refinement question discussed in \cite{Romagnoli2007}.

Common sampling and independence already have substantial precedents.
Huang, Maass, and Ye \cite[Theorem~2.3]{HMY2004} identify the attained
supremum of measure sequence entropies of a fixed partition. More
particularly, their Lemma~3.3 and Theorem~3.5 give conditions on a
measurable cover under which there is \emph{one} sampling sequence with a
strictly positive lower bound for the measure sequence entropies of
\emph{all} its finite measurable refinements. Hence a common positive
threshold over Borel refinements, without further qualification, is not
the contribution asserted here. Their bound uses an invariant measure
and a sequence obtained from that cover; it is not the exact
combinatorial value $\log N(\K)$ for every prescribed sequence with
diverging gaps. The measure-cover quantities and their relation to
independence are developed further by Kerr and Li
\cite[Section~4 of the author version]{KL2009}.

Li \cite[Theorem~2]{Li2002} characterizes topological weak mixing by
sequence-entropy attainment for each regular open cover, with the
sequence chosen for that cover. Huang and Ye
\cite[Theorem~2.1]{HY2009} prove attainment of maximal pattern entropy by
sequence entropy for a fixed open cover, as well as a measurable
counterpart. Independence of finite patterns is part of the established
framework of Kerr and Li \cite{KL2007}. We use the elementary
finite-type form of this independence and give its proof. The additional
step here is to pass from nonmeager Borel atoms to actual orbit patterns
on every prescribed sequence with diverging gaps, and to convert
closed-cover refinement into an exact cardinality optimization. Neither
independence itself nor a new general variational principle is claimed.

There is also a control interpretation. A safe feedback assigns a
finite control block to each state and repeats this assignment at a
fixed period. If all legal blocks have the same endpoint map, feedback
does not alter the state dynamics at block boundaries, but may choose
among overlapping safety domains in any Borel fashion. Theorem~\ref{thm:control}
computes the full-feedback minimax value as $\log r/\tau$, without
restricting feedbacks to continuous or finite-window rules. A positive
example has overlapping safety domains and no nonconstant control-letter
projection forced by every feedback. Finally, we give an exact example
in which feedback changes the successor: the same cover number is three,
while the full-class minimax value is $\log 2$. This is a boundary of the
cover formula, not a counterexample to minimax equality itself.

\section{Counting entropy and the closed-cover theorem}

Let $Q$ be a nonempty compact metric space and let $T:Q\to Q$ be a Borel
map. For a finite Borel partition $\alpha$ of $Q$ and a finite
$F\subset\Nzero$, put
\[
 \Patt_\alpha(F)=
 \left|\left\{(\alpha(T^k x))_{k\in F}:x\in Q\right\}\right|,
\]
where $\alpha(y)$ denotes the atom containing $y$ and empty atoms are
discarded. Thus $\Patt_\alpha(F)$ is the number of nonempty atoms in
$\bigvee_{k\in F}T^{-k}\alpha$. For a strictly increasing sequence
$A=(a_i)_{i\ge1}$ in $\Nzero$, define
\begin{align}
 h_A(T,\alpha)&=\limsup_{n\to\infty}\frac1n
       \log\Patt_\alpha(\{a_1,\ldots,a_n\}),\label{eq:seq}\\
 p_\alpha^*(n)&=\max_{\substack{F\subset\Nzero\\ |F|=n}}
                    \Patt_\alpha(F),
 &\hpat(T,\alpha)&=\lim_{n\to\infty}\frac1n\log p_\alpha^*(n).
 \label{eq:pat}
\end{align}
All logarithms are natural. These are combinatorial counting quantities;
for Borel atoms they are not Shannon entropies of an invariant measure.
The maximum in \eqref{eq:pat} exists because the counts are integers
between $1$ and $|\alpha|^n$. Splitting a window into two parts gives
$p_\alpha^*(n+m)\le p_\alpha^*(n)p_\alpha^*(m)$. Subadditivity proves the
existence of the limit, and $h_A(T,\alpha)\le\hpat(T,\alpha)$.
No invertibility or continuity of the observation is used here.

Let $\K=(K_w)_{w\in W}$ be a finite closed cover of $Q$, and write
\[
 r=N(\K)=\min\{|I|:I\subset W,\ \bigcup_{w\in I}K_w=Q\}.
\]
The set $\mathcal B(\K)$ consists of all finite Borel partitions refining
$\K$: each atom must be contained in at least one $K_w$. Equivalently,
one may first consider all Borel maps $s:Q\to W$ with $x\in K_{s(x)}$,
and then allow arbitrary finite Borel refinements of their fibers.

Let $H$ be a primitive finite $0$--$1$ matrix, with vertex set $V$, and
let
\[
 Y_H=\{y\in V^{\Nzero}:H_{y_jy_{j+1}}=1\text{ for all }j\ge0\}.
\]
Primitivity means that all sufficiently large powers of $H$ have strictly
positive entries. Write $\sigma$ for the one-sided shift.

\begin{theorem}\label{thm:cover}
Suppose $T:Q\to Q$ is topologically conjugate to $\sigma^p:Y_H\to Y_H$
for some primitive $H$ and integer $p\ge1$. Let $\K$ be any finite closed
cover of $Q$. For every strictly increasing $A=(a_i)$ satisfying
\begin{equation}\label{eq:gaps}
 a_{i+1}-a_i\longrightarrow\infty,
\end{equation}
one has
\begin{equation}\label{eq:coverformula}
 \min_{\alpha\in\mathcal B(\K)} h_A(T,\alpha)
 =\min_{\alpha\in\mathcal B(\K)}\hpat(T,\alpha)
 =\log r.
\end{equation}
An ordered difference partition of any minimum subcover attains both
minima, simultaneously for every $A$ satisfying \eqref{eq:gaps}.
Consequently,
\begin{equation}\label{eq:coverminimax}
 \max_A\inf_{\alpha\in\mathcal B(\K)}h_A(T,\alpha)=\log r,
\end{equation}
where the maximum ranges over all strictly increasing infinite sequences
and is attained, for example, by $a_i=(i-1)^2$.
\end{theorem}

The order of quantifiers is part of the assertion. Every $A$ satisfying
\eqref{eq:gaps} works for every allowed partition, although the finite
initial segment discarded in the proof can depend on that partition.
The proof does not interchange an infimum over partitions with a
finite-window limit.

\subsection{A category and concatenation lemma}

We use meagerness relative to $Q$. A Borel set has the Baire property,
so a nonmeager Borel set is comeager in some nonempty open set. Indeed,
the sets differing from open sets by meager sets form a sigma-algebra:
for complements, the additional difference is contained in the
nowhere dense boundary of an open set. A compact metric space is a
Baire space. In particular, a nonempty open set cannot be covered by
countably many nowhere dense sets, as is also seen by choosing nested
nonempty open sets with closures avoiding those sets and diameters
tending to zero.

If $F:Q\to Q$ is continuous and open, preimages of meager sets are
meager. For a closed nowhere dense $N$, the closed set $F^{-1}N$ has
empty interior: otherwise the image of a nonempty open subset would be
a nonempty open subset of $N$. Passing to countable unions proves the
assertion.

\begin{lemma}\label{lem:glue}
Under the dynamical assumptions of Theorem~\ref{thm:cover}, $T$ is
continuous, open, and onto. Given finitely many nonempty cylinder sets
$C_1,\ldots,C_q$ in symbolic coordinates, there is an integer $b\ge1$
such that, for every $k\ge1$, every $t_1,\ldots,t_k\in\{1,\ldots,q\}$,
and all $0\le n_1<\cdots<n_k$ with $n_{j+1}-n_j\ge b$,
\[
 \bigcap_{j=1}^k T^{-n_j}C_{t_j}
\]
is nonempty and open.
\end{lemma}
\begin{proof}
Work on $Y_H$ through the conjugacy. The image under $\sigma$ of a
prefix cylinder of length at least two is the cylinder obtained by
removing its first symbol. The image of a one-symbol cylinder is the
union of the permitted successor cylinders. Hence $\sigma$ is open.
Every vertex has a predecessor, so every infinite path has a predecessor;
thus $\sigma$ is onto. These properties pass to powers and conjugates.

Choose within each $C_j$ a nonempty prefix cylinder, and extend their
defining words, if necessary, to a common length $\ell$. Choose $N$ with
$H^n>0$ for every $n\ge N$, and then choose $b$ with
$pb-\ell+1\ge N$. Place the word for $C_{t_j}$ at coordinate $pn_j$.
Between the end of one placed word and the start of the next there are
$p(n_{j+1}-n_j)-\ell+1$ edges to fill. Primitivity supplies a path of
that length between any two endpoint vertices. Predecessors fill the
finite segment before the first placed word, and successors provide an
infinite continuation after the last. This constructs an actual point
in the intersection. The intersection is open because it is a finite
intersection of preimages of open sets under continuous maps.
\end{proof}

\begin{lemma}\label{lem:borel}
Under the same assumptions, let $\alpha$ be a finite Borel partition
with $q$ nonmeager atoms. For every $A$ satisfying \eqref{eq:gaps} there
exists $m\ge1$ such that
\begin{equation}\label{eq:categorycount}
 \Patt_\alpha(\{a_1,\ldots,a_n\})\ge q^{n-m+1}\qquad(n\ge m).
\end{equation}
In particular $h_A(T,\alpha)\ge\log q$.
\end{lemma}
\begin{proof}
List the nonmeager atoms as $D_1,\ldots,D_q$. The Baire property
provides a nonempty cylinder $C_j$ for which $C_j\setminus D_j$ is
meager. Let $b$ be supplied by Lemma~\ref{lem:glue}, and choose $m$ so
that $a_{i+1}-a_i\ge b$ for all $i\ge m$.

Fix $n\ge m$ and any target labels $t_m,\ldots,t_n$. The set
\[
 G=\bigcap_{i=m}^n T^{-a_i}C_{t_i}
\]
is nonempty and open. Each power $T^{a_i}$ is continuous and open, so
\[
 \bigcup_{i=m}^n T^{-a_i}(C_{t_i}\setminus D_{t_i})
\]
is meager and cannot cover $G$. A point outside this union realizes
all the specified atoms at the sampled times. Distinct target label
strings give distinct patterns; restriction from the full window to
this tail is onto. This proves \eqref{eq:categorycount}. Dividing its
logarithm by $n$ and taking a limit inferior proves the asserted bound.
\end{proof}

\subsection{Proof of the exact formula}

\begin{proof}[Proof of Theorem~\ref{thm:cover}]
Let $\alpha\in\mathcal B(\K)$. Assign to each atom $D$ one index
$w(D)$ with $D\subset K_{w(D)}$. Let $J$ be the set of indices assigned
to the nonmeager atoms. If the closed sets $(K_w)_{w\in J}$ did not
cover $Q$, their complement would be a nonempty open set covered by
the finitely many meager atoms of $\alpha$, which is impossible.
Thus $|J|\ge r$, and the number $q$ of nonmeager atoms is at least $r$.
Lemma~\ref{lem:borel} gives
\[
 h_A(T,\alpha)\ge\log q\ge\log r
\]
for every $A$ satisfying \eqref{eq:gaps}. Closedness of the cover
members is used precisely in this passage from category to a full cover.

Choose a minimum subcover $K_{w_1},\ldots,K_{w_r}$, and set
\begin{equation}\label{eq:ordered}
 E_j=K_{w_j}\setminus\bigcup_{k<j}K_{w_k},\qquad
 \alpha_0=\{E_1,\ldots,E_r\}.
\end{equation}
Every $E_j$ is Borel and is nonempty, since a minimum subcover has no
redundant member. The partition refines $\K$ and has $r$ atoms. Hence
$\Patt_{\alpha_0}(F)\le r^{|F|}$ for all finite $F$, giving
$h_A(T,\alpha_0)\le\hpat(T,\alpha_0)\le\log r$ for every $A$.
The established lower bound and $\hpat\ge h_A$ yield
\eqref{eq:coverformula} and simultaneous attainment.
For arbitrary $A$, the partition $\alpha_0$ bounds the inner infimum
in \eqref{eq:coverminimax} by $\log r$. Any $A$ satisfying
\eqref{eq:gaps} gives the reverse bound and attains the maximum.
\end{proof}

\section{A cover--partition gap at prescribed sparse times}

For any finite cover $\K$ define its combinatorial sequence entropy by
\[
 h_A(T,\K)=\limsup_{n\to\infty}\frac1n
 \log N\left(\bigvee_{i=1}^n T^{-a_i}\K\right).
\]
For a partition this is \eqref{eq:seq}. For overlapping sets, covering
each finite orbit window need not correspond to one fixed refining
partition used at every time.

\begin{corollary}\label{cor:covergap}
Let $q\ge3$, $Q=\{0,\ldots,q-1\}^{\Nzero}$, $T=\sigma$, and
\[
 K_j=\{x:x_0\ne j\}\quad(0\le j<q),\qquad \K=(K_j)_{j=0}^{q-1}.
\]
For every increasing infinite $A$,
\begin{equation}\label{eq:covervalue}
 h_A(\sigma,\K)=\log\frac q{q-1}.
\end{equation}
For every such $A$ with diverging successive gaps,
\begin{equation}\label{eq:refvalue}
 \min_{\alpha\in\mathcal B(\K)}h_A(\sigma,\alpha)=\log2
 >h_A(\sigma,\K).
\end{equation}
\end{corollary}
\begin{proof}
The cover has no single member equal to $Q$, and any two distinct
members cover $Q$. Thus $r=2$, and \eqref{eq:refvalue} follows from
Theorem~\ref{thm:cover} once \eqref{eq:covervalue} is proved.

For a window $F$ of size $n$, the full shift has all $q^n$ possible
input vectors. Each member of the joined cover specifies one excluded
letter at each coordinate, and covers $(q-1)^n$ of these vectors.
Consequently the minimum number $c(n)$ of such boxes is at least
$(q/(q-1))^n$. For completeness, the following elementary probabilistic
covering bound gives the matching exponent. Choose $R$ excluded-letter
vectors independently and uniformly from the $q^n$ possibilities. A
fixed input vector is left uncovered with probability at most
\[
 \left(1-\left(\frac{q-1}{q}\right)^n\right)^R
 \le \exp\left(-R\left(\frac{q-1}{q}\right)^n\right).
\]
Taking
$R=\left\lceil(q/(q-1))^n(n\log q+1)\right\rceil$ makes the sum of
these probabilities over all input vectors smaller than one. Some
choice therefore covers every vector. We have proved
\[
 \left(\frac q{q-1}\right)^n\le c(n)
 \le\left\lceil\left(\frac q{q-1}\right)^n(n\log q+1)\right\rceil.
\]
The count depends only on $n$, not on the positions of the sampled
coordinates, so its exponential rate proves \eqref{eq:covervalue}.
\end{proof}

The excluded-letter cover and its covering exponent are already present
in Romagnoli \cite{Romagnoli2007}; the calculation above is included to
identify the sparse-time quantity without an external dependency.
The consequence of Theorem~\ref{thm:cover} is the exact minimum over
\emph{all Borel} refinements for a sequence fixed in advance. It does not
assert that every Borel refinement has consecutive-time entropy at least
$\log2$, or that every sampling sequence gives that lower bound.

\section{Fixed-period safe feedbacks with a common successor}

Let $X$ be a compact metric space, $U$ a nonempty finite control set, and
$f_u:X\to X$ continuous for every $u\in U$. Let $Q\subset X$ be a
nonempty compact controlled-invariant set: at each $x\in Q$ at least
one control keeps the next state in $Q$. All infinite control strings
are allowed. Fix an integer period $\tau\ge1$ once and for all, and put
$W=U^\tau$. For $w=(u_0,\ldots,u_{\tau-1})$ write
\[
 \varphi(0,x,w)=x,\qquad
 \varphi(j,x,w)=f_{u_{j-1}}\circ\cdots\circ f_{u_0}(x),
\]
and define
\begin{equation}\label{eq:safe}
 K_w=\{x\in Q:\varphi(j,x,w)\in Q\ (0\le j\le\tau)\},
 \qquad g_w(x)=\varphi(\tau,x,w).
\end{equation}
The $K_w$ are closed and cover $Q$. The coverage follows by choosing
legal single-step controls successively for the fixed finite horizon.

The class $\Ssafe_\tau$ consists of \emph{all} Borel selectors $s:Q\to W$
such that $x\in K_{s(x)}$. Its closed-loop block map is
$T_s(x)=g_{s(x)}(x)$. All Borel mixtures of legal selectors belong to
$\Ssafe_\tau$. No continuity or dependence on finitely many coordinates
is imposed. A selector is used without an additional clock, register,
or feedback memory.

For a finite window $F$, let
\[
 \Patt_s(F)=\left|\{(s(T_s^k x))_{k\in F}:x\in Q\}\right|.
\]
Define $h_A(s)$ and $\hpat(s)$ from these counts as in
\eqref{eq:seq}--\eqref{eq:pat}, but dividing by $n\tau$ rather than
$n$. Thus normalization is per sampled control block, divided by its
fixed length; it is \emph{not} division by elapsed time $\tau a_n$.
Put
\begin{equation}\label{eq:LM}
 L_\tau=\sup_A\inf_{s\in\Ssafe_\tau}h_A(s),\qquad
 M_\tau=\inf_{s\in\Ssafe_\tau}\hpat(s).
\end{equation}
Always $L_\tau\le M_\tau$.

The counts use actual closed-loop trajectories. If instead one takes
the closure of their one-sided name set, the finite patterns are
unchanged: a cylinder specifying finitely many coordinates is open, so
if it meets that closure, it meets the original name set. Taking a name
closure therefore introduces no extra finite-window witnesses.

\begin{lemma}\label{lem:merge}
The infima in \eqref{eq:LM}, and the inner infimum at any fixed $A$, are
unchanged if one allows every finite Borel safe partition with a
control block assigned to each atom, and counts atom names rather than
control-block names.
\end{lemma}
\begin{proof}
Merge all atoms assigned the same block. Their finite union is Borel
and remains inside that block's safety domain. The resulting selector
has exactly the same closed-loop state map. The map from original atom
names to block names is coordinatewise and onto the actual block-name
patterns, so every block-pattern count is no larger than the original
atom-pattern count. Conversely, the nonempty fibers of any legal
selector form a legal finite Borel safe partition whose atom and
block counts agree. Both directions of the infimum equality follow.
\end{proof}

\begin{theorem}\label{thm:control}
Assume in addition that there is a map $T:Q\to Q$ such that
\begin{equation}\label{eq:common}
 g_w(x)=T(x)\qquad\text{for every }w\in W\text{ and every }x\in K_w,
\end{equation}
and that $T$ is conjugate to $\sigma^p$ on a mixing one-sided shift of
finite type. With $r=N((K_w)_{w\in W})$, for every $A$ satisfying
\eqref{eq:gaps} one has
\begin{equation}\label{eq:controlformula}
 \min_{s\in\Ssafe_\tau}h_A(s)
 =\min_{s\in\Ssafe_\tau}\hpat(s)
 =L_\tau=M_\tau=\frac{\log r}{\tau}.
\end{equation}
The maximum defining $L_\tau$ is attained by every such $A$, and an
ordered minimum-subcover selector attains all the indicated minima.
The same assertions hold over all finite Borel safe partitions and
their arbitrary finite Borel refinements.
\end{theorem}
\begin{proof}
For every legal selector, \eqref{eq:common} gives $T_s=T$ pointwise.
Its fibers form a finite Borel refinement of $(K_w)_{w\in W}$.
Theorem~\ref{thm:cover}, with normalization by $\tau$, gives the lower
bound at every $A$ with diverging gaps. If $w_1,\ldots,w_r$ form a
minimum subcover, define $s_0(x)=w_j$ on the sets $E_j$ of
\eqref{eq:ordered}. This is Borel and legal at every point. Each selected
block stays in $Q$ at its intermediate times by \eqref{eq:safe}, and
the endpoint again lies in $Q$, so recursion gives an infinite safe
closed-loop orbit for every initial state. The selector uses only $r$
blocks and thus gives the upper bound for every window and every
sampling sequence. This proves attainment, \eqref{eq:controlformula},
and the outer maximum. Lemma~\ref{lem:merge} proves the last assertion.
\end{proof}

The standard safe-partition framework is described by Kawan
\cite[Chapter~2]{Kawan2013}. His minimization over control periods in
invariance entropy is a different optimization. Here $\tau$ is fixed,
intermediate safety is retained, and the exact value comes from the
additional common-successor and finite-type hypotheses.

\subsection{A positive overlapping example at period two}

Take
\[
 Q=\{0,1,2\}^{\Nzero},\quad X=Q\sqcup\{\dagger\},\quad
 U=\{0,1,2\},\quad\tau=2,
\]
where the isolated point $\dagger$ lies outside $Q$. Define
\[
 f_u(x)=
 \begin{cases}\sigma x,&x\in Q,\ x_0\ne u,\\
 \dagger,&x\in Q,\ x_0=u,\end{cases}
 \qquad f_u(\dagger)=\dagger.
\]
Each $f_u$ is continuous since the cases are clopen. For $w=(a,b)$,
\[
 K_{(a,b)}=\{x:x_0\ne a,\ x_1\ne b\},\qquad
 g_{(a,b)}|_{K_{(a,b)}}=\sigma^2.
\]
All nine control blocks are available to Borel feedbacks, and each
state permits four of them. A safety domain covers four of the nine
two-symbol input blocks, so two domains cannot cover $Q$. The three
diagonal blocks $(0,0),(1,1),(2,2)$ do cover $Q$, since some symbol
differs from both $x_0$ and $x_1$. Hence $r=3$, and
\[
 L_2=M_2=\frac{\log3}{2}.
\]
An attaining selector is
\[
 c(x)=\min(\{0,1,2\}\setminus\{x_0,x_1\}),\qquad
 s_0(x)=(c(x),c(x)).
\]
It is clopen and legal everywhere. Its actual block-boundary states
are $\sigma^{2k}x$. To prescribe any sequence of diagonal labels, use
successive input blocks $11$, $00$, and $01$ for the desired labels
$0$, $1$, and $2$, respectively. Therefore its actual name set is the
full shift on these three blocks. The lower bound, however, holds for
every legal Borel rule using any of the nine blocks.

This example does not preassign a nonconstant common control-letter
observation to every feedback. More precisely, suppose a map
$\rho:U^2\to V'$ and an observation $F:Q\to V'$ satisfy
$\rho(s(x))=F(x)$ for every legal Borel selector and every $x$.
Any two blocks $(a,b)$ and $(a',b')$ are colegal at a state beginning
with $ij$, where $i\notin\{a,a'\}$ and $j\notin\{b,b'\}$.
A legal baseline selector may be changed at this single state to use
either block, while remaining Borel and legal. Hence their $\rho$
values agree. All pairs of blocks are colegal somewhere, so $\rho$
is constant. This rules out this particular forced-letter explanation;
it makes no claim about every possible dynamical factor.

\section{A feedback-dependent boundary with an exact full-class value}

The common-successor assumption cannot be removed solely because each
branch is continuous and open. We give a system with individual
prefix cylinders as safety domains and compute its full-class value.

Let $Q=\{0,1\}^{\Nzero}$, let $\mathcal C=\{0,10,11\}$ be a set of
finite words, and take $U=\mathcal C\times\{0,1\}$ and $\tau=1$.
For $e\in\{0,1\}$ let $P_e$ add $e$ modulo two to every coordinate.
On $X=Q\sqcup\{\dagger\}$ define
\begin{equation}\label{eq:boundarysystem}
 f_{(c,e)}(x)=
 \begin{cases}P_e\sigma x,&x\in[c],\\
 \dagger,&x\in Q\setminus[c],\end{cases}
 \qquad f_{(c,e)}(\dagger)=\dagger.
\end{equation}
The cylinders $[0],[10],[11]$ are disjoint and cover $Q$. Each state
has exactly two legal controls with different successors. The global
maps $f_{(c,e)}$ are continuous, and their legal restrictions are open.
The safety domains are $K_{(c,e)}=[c]$, and $r=3$.

\begin{proposition}\label{prop:boundary}
For the system \eqref{eq:boundarysystem}, over all legal Borel feedbacks
and arbitrary Borel mixtures,
\begin{equation}\label{eq:boundaryvalue}
 \min_s h_{A_0}(s)=\min_s\hpat(s)=
 \max_A\inf_s h_A(s)=\log2<\log r,
 \qquad A_0=(0,1,2,\ldots).
\end{equation}
There is one clopen legal feedback $s_{\mathrm f}$ satisfying
$h_A(s_{\mathrm f})=\log2$ for every increasing infinite $A$.
The same full-class values hold if arbitrary finite Borel atom
refinements are allowed.
\end{proposition}
\begin{proof}
Let $c(x)$ be the unique member of $\mathcal C$ with $x\in[c(x)]$,
and set
\[
 s_{\mathrm f}(x)=(c(x),x_1),\qquad
 T_{\mathrm f}x=P_{x_1}\sigma x.
\]
This is a clopen selector, legal at every state. Its four nonempty
initial control fibers, for labels $(0,0),(0,1),(10,0),(11,1)$, are
$[00],[01],[10],[11]$, respectively. For all $n\ge1$, direct induction
gives
\begin{equation}\label{eq:fliporbit}
 T_{\mathrm f}^n x=P_{x_n}\sigma^n x,
 \qquad s_{\mathrm f}(T_{\mathrm f}^n x)
        =(0,x_{n+1}\mathbin\oplus x_n),
\end{equation}
where $\oplus$ is addition modulo two. In the induction step the
chosen flip is $x_{n+1}\oplus x_n$; it cancels the previous flip and
leaves $P_{x_{n+1}}\sigma^{n+1}x$. The first coordinate is always zero
after time zero, which also proves all-time safety.

For every finite nonempty window $F$, with $n=|F|$, one has
\begin{equation}\label{eq:flipcount}
 \Patt_{s_{\mathrm f}}(F)=
 \begin{cases}2^n,&0\notin F,\\2^{n+1},&0\in F.\end{cases}
\end{equation}
The four initial labels and two later labels give the upper bounds.
To prove the lower bounds, prescribe any of the four initial labels
if $0\in F$, thereby fixing $(x_0,x_1)$; if $0\notin F$, set
$x_0=x_1=0$. At positive sampled times prescribe the desired second
components $d_k$, and put $d_k=0$ at the other positive times.
Recursively set $x_{k+1}=x_k\oplus d_k$ for every $k\ge1$.
This constructs a full state in $Q$ whose actual trajectory under the
same selector has the required sampled controls by
\eqref{eq:fliporbit}. Thus all counted patterns are realized by true
infinite closed-loop trajectories. Equation~\eqref{eq:flipcount}
implies $h_A(s_{\mathrm f})=\hpat(s_{\mathrm f})=\log2$ for every $A$.

For the opposite bound fix any legal Borel selector $s$. If its first
$n$ actual controls from $x$ are $(c_k,e_k)$, let
$E_0=0$ and $E_k=e_0\oplus\cdots\oplus e_{k-1}$. Since shifts commute
with flips,
\[
 T_s^k x=P_{E_k}\sigma^k x.
\]
Let $b(c)$ be the first bit of the word $c$. Safety at time $k$ gives
\begin{equation}\label{eq:decode}
 x_k=b(c_k)\oplus E_k\qquad(0\le k<n).
\end{equation}
Thus one actual length-$n$ control word determines at most one input
prefix of length $n$. Every binary input prefix has an extension in
$Q$, and the feedback is defined and safe on every such extension.
Distinct prefixes must yield distinct actual control words. Hence
\[
 \Patt_s(\{0,\ldots,n-1\})\ge2^n
\]
for every $n$ and every legal $s$. This proves $h_{A_0}(s)\ge\log2$
and $\hpat(s)\ge\log2$. Combining these inequalities with
$s_{\mathrm f}$ proves \eqref{eq:boundaryvalue}, including attainment.
Lemma~\ref{lem:merge} handles atom refinements.
\end{proof}

The obstruction has a precise orbit interpretation. In this example
$T_{\mathrm f}(Q)=[0]$: for the reverse inclusion, given $y\in[0]$,
take $x=0y$. Therefore neither $[10]$ nor $[11]$ is visited at a
positive time. The initial fibers with those labels are nonmeager,
but their labels cannot be freely concatenated in the future.
The map $T_{\mathrm f}$ is itself continuous and open, as follows by
restricting to the four clopen two-symbol cylinders. It is not onto
$Q$.

Even preservation of meager preimages does not repair this problem.
For any selector $s$ and any $n$, the value of $T_s^n x$ is either
$\sigma^n x$ or $P_1\sigma^n x$. Hence, for a meager set $N$,
\[
 (T_s^n)^{-1}N\subset
 (\sigma^n)^{-1}N\ \cup\ (P_1\sigma^n)^{-1}N
\]
is meager. What fails is the nonempty open intersection used in
Lemma~\ref{lem:glue} for the \emph{selected} dynamics. Also, the
same $s_{\mathrm f}$ gives
\[
 \inf_s h_A(s)\le\log2=\log r-\log(3/2)
 \quad\text{for every }A.
\]
The failure of the cover formula is therefore a uniform numerical
gap over all samplings, rather than a failure at one finite window.
It does not refute $L_\tau=M_\tau$: equality holds in this very example.
Nor does the proof determine the inner infimum for every other fixed
sampling sequence.

\section{Scope of the conclusions}

The exact closed-cover formula combines three specific features:
closed safety constraints force enough nonmeager labels, an open common
map preserves category under preimages, and finite-type concatenation
realizes every finite choice of those labels at sufficiently separated
times. This makes a prescribed divergent-gap sequence work for an
uncountable family of Borel observations without a countable
approximation argument. The ordered subcover gives an actual minimizer,
so no assumption about attainment of an abstract infimum is needed.

The control application retains every legal Borel selector, but has a
substantial dynamical restriction: legal blocks share their endpoint.
The six-control example identifies a failure of extending the cover
number formula to changing successors. Its exact value uses an explicit
normalizing feedback and input-prefix decoding; it does not supply a
theory for arbitrary feedback-dependent branches. The equality
$L_\tau=M_\tau$ for general finite-control, fixed-period systems over
the full Borel class remains undecided by these results. No conclusion
about simultaneous attainment of each feedback's individual maximal
pattern entropy follows from the common minimum threshold proved here.

\begin{thebibliography}{99}

\bibitem{HMY2004}
Huang, W., Maass, A., \& Ye, X. (2004).
Sequence entropy pairs and complexity pairs for a measure.
\emph{Annales de l'Institut Fourier, 54}(4), 1005--1028.
\doi{10.5802/aif.2041}.

\bibitem{HY2009}
Huang, W., \& Ye, X. (2009).
Combinatorial lemmas and applications to dynamics.
\emph{Advances in Mathematics, 220}(6), 1689--1716.
\doi{10.1016/j.aim.2008.11.009}.

\bibitem{Kawan2013}
Kawan, C. (2013).
\emph{Invariance entropy for deterministic control systems: An introduction}
(Lecture Notes in Mathematics, Vol.~2089). Springer.
\doi{10.1007/978-3-319-01288-9}.

\bibitem{KL2007}
Kerr, D., \& Li, H. (2007).
Independence in topological and C*-dynamics.
\emph{Mathematische Annalen, 338}, 869--926.
\doi{10.1007/s00208-007-0097-z}.
Author version: \href{https://arxiv.org/abs/math/0603585v4}{arXiv:math/0603585v4}.

\bibitem{KL2009}
Kerr, D., \& Li, H. (2009).
Combinatorial independence in measurable dynamics.
\emph{Journal of Functional Analysis, 256}(5), 1341--1386.
\doi{10.1016/j.jfa.2008.12.014}.
Author version: \href{https://arxiv.org/abs/0705.3424v1}{arXiv:0705.3424v1}.

\bibitem{Li2002}
Li, S. (2002).
Topological sequence entropy and topologically weak mixing.
\emph{Bulletin of the Australian Mathematical Society, 66}(2), 207--211.
\doi{10.1017/S0004972700040053}.

\bibitem{Romagnoli2007}
Romagnoli, P.-P. (2007).
A remark on the topological entropies of covers and partitions.
\emph{Studia Mathematica, 182}(3), 273--281.
\doi{10.4064/sm182-3-6}.

\end{thebibliography}
\end{document}
